\cotefile{156-8}
\scanpage{34}
\mspage{25}
% archive sheet number: 34
% inv: text | top | Magasins triviaux
\uline{Magasins triviaux}\quad Def : Un magasin est dit ``trivial'' si $X\overset{\circ}{\ll}Y\Longleftrightarrow X=Y$, et si $X\mathrel{|\circ|}Y\Longleftrightarrow X\ne Y$, donc la seule donnée qui reste à déterminer est $X\leq Y$. Mais on \unsure{vérifie} \inserted{trivialement} que pour \deleted{tout} ens. ordonné $(M,\leq)$, si on prend \deleted{$X\overset{\circ}{\ll}Y$}
% inv: display | upper third | X circle-ll Y
\[
X\overset{\circ}{\ll}Y\overset{\text{déf}}{\Longleftrightarrow}X=Y,\qquad X\mathrel{|\circ|}Y\overset{\text{déf}}{\Longleftrightarrow}X\ne Y,
\]
% inv: text | upper middle | on trouve une structure
on trouve une structure de magasin.

% inv: deletion | upper middle | Ce magasin ne satisfait
\deleted{Ce magasin ne satisfait Mag L 1 que si \unsure{tout}} \deleted{\unsure{fig.}} \deleted{$L=$ ens des \unsure{éléts} minimaux}

% inv: display | middle | X ll Y iff X le Y
\[
\begin{aligned}
X\ll Y&\Longleftrightarrow X\leq Y\\
X\unsure{\gtrless|\bowtie}Y&\qquad\text{tjrs satisfait}\\
X\parallel Y&\qquad\text{\unsure{,,,}}\quad\widetilde X\cap\widetilde Y=\varnothing\quad\text{i.e.}\quad\{X,Y\}\text{ pas minoré}.
\end{aligned}
\]
% inv: display | middle | L equals ens des éléments minimaux
\[
L=\text{ens des éléments minimaux de }X
\]
% inv: display | middle lower | omb(X)
\[
\begin{aligned}
\mathrm{omb}(X)&=\{a\in L\mid a\leq X\}\\
\mathrm{omb}(X)^\circ&=\deleted{\{a\in L\mid a<X\}}\begin{cases}\varnothing\text{ si }X\notin L\\\mathrm{omb}(X)=\{X\}\text{ si }X\in L\end{cases}
\end{aligned}
\]

% inv: margin note | slanted left upper | NB Conditions équivalentes
% margin note written sideways
% slanted
\margin{\uline{NB} Conditions \unsure{équivalentes} pour un magasin $M$ :\par
(i) $M$ trivial\par
(ii) $X\gtrless Y$ tjrs satisfait\par
(iii) $\unsure{\mathrm{FigMag}|\mathrm{Fig}_{\mathrm{Mag}}}(M)=$ ens. des figures \unsure{initiales|filtrantes} pour $\leq$\par
(iv) $M\in\mathrm{Fig}_{\mathrm{mag}}(M)$\par
(v) $X\ne Y\Longrightarrow X\mathrel{|\circ|}Y$}
% inv: margin note | slanted left middle | Ces conditions impliquent
% margin note written sideways
% slanted
\margin{Ces conditions impliquent les conditions suivantes qui \unsure{sont} \unsure{équivalentes}, mais qui n'impliquent pas les précédentes\par
(vi) $X\overset{\circ}{\ll}Y\Longrightarrow X=Y$\par
(vii) $X\ll Y\Longrightarrow X\leq\deleted{\unsure{\gtrless}}Y$}

% inv: text | lower middle | La condition Mag L 1
La condition \uuline{Mag L 1} i.e. $\mathrm{omb}(X)^\circ\ne\unsure{0|\varnothing}\ \forall X$, \unsure{signifie} que tout $X$ est minimal i.e. la relation d'ordre est discrète. On voit qu'un magasin trivial ne satisfait \unsure{quasiment|pratiquement} jamais Mag L 1, et encore moins Mag L 2 --- sauf s'il est doublement trivial $(X\leq Y\Longleftrightarrow X=Y)$.

% inv: margin note | slanted left lower | Les magasins triviaux
% margin note written sideways
% slanted
\margin{Les magasins triviaux jouent le rôle d'ensembles d'indices pour des familles de magasins \unsure{généraux}, ou des \uline{maquettes} \unsure{combinatoires} de figures. On peut \unsure{appeler} \unsure{simplement} un tel magasin une \uline{maquette} (!), (une \unsure{m.} \unsure{squelette} (\unsure{catégorique})).}

% inv: text | bottom | Figure indexée
\uline{Figure indexée} par un ens. ordonné $I$ : figure $F$, \unsure{munie} d'un isom.
% inv: diagram | bottom | I isomorphic to tilde F
\diagram{\begin{tikzcd}[ampersand replacement=\&]
\deleted{\unsure{\in}}\quad I \arrow[r,"\sim",dash] \& \widetilde F\\
i \arrow[r,mapsto] \& X_i
\end{tikzcd}}

\scanpage{35}
\mspage{26}
% archive sheet number: 35
% inv: text | top | Ainsi une figure indexée
Ainsi, une figure indexée \inserted{par $I$} définit une application
% inv: display | top | varphi I to M
\[
\varphi:I\longrightarrow M
\]
% inv: text | upper third | qui est croissante
qui est \uline{croissante} pour $\leq$ sur $M$, \deleted{Quelles} et $F$ est déterminée par $\varphi$ \deleted{\unsure{comme}} avec
% inv: display | upper third | tilde F equals varphi(I)
\[
\widetilde F=\varphi(I).
\]
% inv: text | upper middle | Quelles conditions
Quelles conditions sur $\varphi$, pour que ça corresponde bien à une figure indexée ? Il faut

% inv: text | upper middle | a) varphi injectif
 a) $\varphi$ injectif, et \deleted{$=$} \inserted{induit} un isom. d'ens. ordonnés $I\simeq\varphi(I)$

% inv: text | middle | et b) varphi(I) est une figure
et\qquad b) $\varphi(I)$ est une figure. \deleted{\unsure{Je dis que}} \deleted{Cela signifie} \deleted{deux choses}

% inv: display | middle | 1 degrees forall i j
\[
1^\circ)\quad\forall i,j\in I,\quad X_i\unsure{\gtrless|\bowtie}X_j\quad\deleted{\unsure{\text{i.e.}}}\quad\text{et}
\]
% inv: text | middle | 2 degrees varphi(I)
$2^\circ)$\quad $\varphi(I)$ est une partie fermée de $M$.

% inv: text | middle lower | La première condition
La première condition revient à celle-ci ; \deleted{qui implique aussi \unsure{a)}}
% inv: display | middle lower | b) restriction
\[
b)\quad\forall i\in I,\quad\varphi|I_{\leq i}:I_{\leq i}\xrightarrow{\sim}\widetilde X_i
\]
% inv: display | middle lower | c) Si i not equal j
\[
c)\quad\text{Si }i\ne j,\quad X_i\mathrel{|\circ|}X_j
\]

% inv: text | lower third | En d'autres termes
En d'autres termes, \unsure{je} dis $X\longrightarrow X_i$,

% inv: text | lower third | Proposition
\uline{Proposition}\quad Pour que $\varphi:I\longrightarrow M$, définisse une figure indexée par $I$, il f et s que

% inv: text | lower third | a) varphi croissante
 a) $\varphi$ \deleted{injectif} croissante (ds $M,\leq$) et \uline{injective}.

% inv: text | lower third | b) application induite
 b) $\forall i\in I$, l'application induite
% inv: display | bottom third | tilde i equals I le i
\[
\deleted{\unsure{\mathrm{de}}}\quad\widetilde i=I_{\leq i}\xrightarrow{\sim}\widetilde X_i=M_{\leq X_i}
\]
% inv: text | bottom third | soit bijective
soit \deleted{\unsure{bijective} (\unsure{compatible} \unsure{toujours|tjrs} \unsure{aux} \unsure{relations} \unsure{d'ordre}} $=$ un isom. d'ensembles ordonnés)

% inv: text | bottom | c) Si i not equal j
 c) Si $i\ne j$ (i.e $i\mathrel{|\circ|}j$ pour la structure de magasin trivial sur $I$ !) \unsure{:} $X_i\mathrel{|\circ|}X_j$.

\scanpage{36}
\mspage{27}
% archive sheet number: 36
% inv: text | top | En effet b) implique
En effet, b) implique déjà que $\varphi(I)$ est fermée. Il reste à prouver que \deleted{\unsure{les} \unsure{figures} $X_i$} pour $i,j\in I$, on a $X_i\unsure{\gtrless|\bowtie}X_j$. Donc que
% inv: display | upper third | X prime in tilde Xi setminus
\[
X'\in\widetilde X_i\setminus\widetilde X_i\cap\widetilde X_j,\quad Y'\in\widetilde X_j\setminus\widetilde X_i\cap\widetilde X_j\Longrightarrow X'\mathrel{|\circ|}Y'.
\]
% inv: text | upper third | Or par b)
Or par \deleted{\unsure{a)}} b) on aura $X'=X_{i'}$, avec $i'\leq i$, idem pour $Y'=\deleted{\unsure{Y}}X_{j'}$, avec $j'\leq j$ ; d'ailleurs comme $X'\ne Y'$ on aura $i'\ne j'$ par a), Donc on a $X_{i'}\mathrel{|\circ|}X_{j'}$ par c), qed.

% inv: text | middle | Variante polyédrale
\deleted{Morphismes de \unsure{mag}} \inserted{Variante polyédrale :} On suppose maintenant l'ens. ordonné $I$ polyédral. On veut se borner aux figures polyédrales. Cela \deleted{\unsure{est une}} \unsure{se fera} sans frais ; \deleted{\unsure{tout}} \unsure{ceci est|sera} automatique, puisque la condition polyédrale sur une figure se lit sur sa maquette \deleted{$\unsure{\widetilde E}$} $\widetilde F$, et signifie justement que celle-ci est un ens. ordonné polyédral.

% inv: text | lower middle | Morphismes de maquettes
\uline{Morphismes de maquettes}\quad $\varphi:M\longrightarrow M'$ : application telle que si $X,Y\in M$, $X'=\varphi(X)$, $Y'=\varphi(Y)$, on ait
% inv: display | lower third | a) b) c) morphism conditions
\[
\begin{aligned}
a)\quad X\leq Y&\Longrightarrow X'\leq Y'\\
b)\quad X\overset{\circ}{\ll}Y&\Longrightarrow X'\overset{\circ}{\ll}Y'\\
c)\quad X\mathrel{|\circ|}Y&\Longrightarrow X'\mathrel{|\circ|}Y'
\end{aligned}
\]
% inv: insertion | after c) | plus d) plus loin page suiv
\inserted{[+ \textcircled{d} plus loin page suiv.]}
% inv: margin note | right beside a) b) | varphi croissante
\margin{$\varphi$ croissante pour $\leq$, $\overset{\circ}{\ll}$ (alors aussi pour $\ll$)}
% inv: text | lower third | on en conclut
on en conclut
% inv: display | lower third | X ll Y implies
\[
X\ll Y\Longrightarrow X'\ll Y'
\]
% inv: text | bottom | Mais aura-t-on aussi
Mais aura-t-on aussi
% inv: display | bottom | questioned implications
\[
\begin{aligned}
X\gtrless Y&\overset{?}{\Longrightarrow}X'\gtrless Y'\\
X\parallel Y&\overset{?}{\Longrightarrow}X'\gtrless Y'
\end{aligned}
\]
% inv: margin note | slanted left lower | Peut-être vaut-il mieux
% margin note written sideways
% slanted
\margin{Peut-être vaut-il mieux parler de \unsure{comorphismes} ? Ça \unsure{ressemble} plus aux propriétés d'une application ``\unsure{image|images} \unsure{inverse|inverses}'' ! Mais alors, ne serait-il pas plus naturel de la regarder tout de suite comme $\widetilde F\longrightarrow\widetilde F'$ des \unsure{étalages|étalées}, sans supposer \unsure{fixé un} $M\longrightarrow M'$ ??}
% inv: margin note | right lower | Est-on trop exigeant
\margin{(Est-on trop exigeant ??) Ça semble raisonnable dans le cas d'un \uuline{plongement} des maquettes \unsure{surtout|seulement} $\rightarrow$. On peut songer à demander \deleted{\unsure{aussi}} \inserted{seulement} $X\gtrless Y\deleted{\Longleftarrow}\Longrightarrow X'\gtrless Y'$ ce qui est moins extravagant\ldots}
% inv: margin note | bottom right | NB Si M est un magasin trivial
\margin{\uline{NB} Si $M$ est un magasin trivial, ces conditions impliquent $\varphi$ injective.}

\scanpage{37}
\mspage{28}
% archive sheet number: 37
% inv: deletion | top | Il semble qu'il manque
\deleted{Il \unsure{semble} qu'il manque une condition pour pouvoir le conclure !}

% inv: text | top | Il semble qu'on ait besoin
Il semble qu'on ait besoin aussi de la condition
% inv: text | upper third | d) forall X in M
 d) $\forall X\in M$, $\varphi$ induit une \uline{bijection}
% inv: display | upper third | tilde X to tilde X prime
\[
\widetilde X\xrightarrow{\sim}\widetilde X'\qquad\text{où}\quad X'=\varphi(X)
\]
% inv: margin note | slanted left upper | dated condition d)
% margin note written sideways
% slanted
\margin{\unsure{12|11}.7.86\par Cette condition d) devrait \unsure{rester|être} \unsure{facultative}. Il y a des ``\unsure{sous-mag.}'' qui ne sont pas fermés par $\unsure{\leq|\Delta}$ !}

% inv: text | upper middle | Corollaire 1
\uline{Corollaire 1} Les conditions a), c), d) impliquent que $\forall F\in\mathrm{Fig}(M)$, on a $\varphi(F)\in\mathrm{Fig}(M')$, et $\varphi|F$ est injective \inserted{et $\simeq$ d'ens. ordonnés}.

% inv: text | middle | Immédiat par la proposition
Immédiat par la proposition précédente \inserted{(si on pose $I=\widetilde F$)} si $\varphi$ est injective. \deleted{Dans le cas général,} Il suffit de vérifier que $\varphi|F$ est injective. Mais si $X,Y\in F$, \deleted{\unsure{si}} $X\ne Y$, on aura $X\mathrel{|\circ|}Y$, donc $\varphi(X)\mathrel{|\circ|}\varphi(Y)$ par c), donc $\varphi(X)\ne\varphi(Y)$ OK ! \deleted{\unsure{Montrons que} $\varphi$} \inserted{C'est la prop qui} montre de plus que $\varphi$ est un isom. ordonné\ldots

% inv: text | middle | Corollaire 2
\uline{Corollaire 2}\quad\deleted{Si $X\unsure{\gtrless|\bowtie}Y$} Si $X'=\varphi(X)$, $Y'=\varphi(\unsure{Y})$
% inv: display | middle | compatibility and parallel preserved
\[
\begin{aligned}
X\gtrless Y&\Longrightarrow X'\gtrless Y'\\
X\parallel Y&\Longrightarrow X'\parallel Y'.
\end{aligned}
\]
% inv: text | middle lower | En effet X relation Y
En effet, $X\gtrless Y$ signifie que $X,Y$ font partie d'une figure \inserted{$F$}, et on applique Cor. 1. Plus \deleted{Si $X\parallel Y$, \unsure{on a} $X'$ \unsure{on aura} $\mathrm{Inf}(X,Y)$} précisément

% inv: text | lower third | Cor 3
\uline{Cor 3}\quad Soit $F$ figure dans $M$. Alors pour $X,Y\in\widetilde F$, et $X'=\varphi(X)$, $Y'=\varphi(Y)$, on a
% inv: display | lower third | deleted relations then two equivalences
\[
\begin{aligned}
&\deleted{X\gtrless Y\Longleftrightarrow X'\gtrless Y'}\\
&\deleted{X\parallel Y}\\
&X\leq Y\Longleftrightarrow X'\leq Y'\\
&\deleted{\unsure{X\gtrless Y}}\\
&X\parallel Y\Longleftrightarrow X'\parallel Y'
\end{aligned}
\]
% inv: margin note | right lower | se généralise
\margin{\unsure{se généralise} aux sous-figures quelconques de $F$}

% inv: text | bottom | La première
La première\inserted{$\unsure{*}$} \unsure{équivalence} \deleted{\unsure{provient}} déjà \unsure{vue|vu} (Cor. 1). La deuxième provient du fait que $X\parallel Y$ se lit dans l'ens ordonné $I=\widetilde F$ ($\{X,Y\}$ non minoré) lequel est isom. à $I'=\varphi(I)$.

\scanpage{38}
\mspage{29}
% archive sheet number: 38
% inv: text | top | Cor.4 Soient F G
\uline{Cor. 4}\quad Soient $F,G\in\mathrm{Fig}(M)$, avec $F\unsure{\gtrless|\bowtie}G$. \deleted{Alors} \unsure{Soient|Soit} $F'=\varphi(F)$ \deleted{et} $G'=\varphi(G)$, alors
% inv: display | upper third | union and intersection
\[
F'\gtrless G',\quad\varphi(F\cup G)=\varphi(F)\cup\varphi(G),\quad\varphi(F\cap G)=\varphi(F)\cap\varphi(G)
\]
% inv: text | upper third | Dém On applique
\uline{Dém}\quad On applique cor 1 à $F\cup G$.

% inv: text | upper middle | Cor.5 Si F ll G
\uline{Cor. 5}\quad Si $F\ll G$, alors $\underbrace{\varphi(F)}_{F'}\ll\underbrace{\varphi(G)}_{G'}$, et on a commutativité de
% inv: diagram | middle | square tilde F tilde G
\diagram{\begin{tikzcd}[ampersand replacement=\&]
\widetilde F \arrow[r] \arrow[d,"\varphi"'] \& \widetilde G \arrow[d,"\varphi"]\\
\widetilde F' \arrow[r] \& \widetilde G'
\end{tikzcd}}

% inv: text | middle | Cor.6 Compatibilité
\uline{Cor. 6}\quad Compatibilité de $\varphi$ avec l'opération de raffinement induit sur les figures.

% inv: text | lower middle | J'ai pas l'impression
J'ai pas l'impression d'avoir bien saisi le rôle exact de l'hypothèse c), qui semble très forte. Dans le cas d'une opération d'image inverse par une application ensembliste, elle semble supposer que cette dernière est \uline{surjective}. Mais dans ce cas, $\varphi$ \deleted{est} \inserted{serait} injective, on pourrait aussi la supprimer, \deleted{\unsure{ceci à supprimer}}

% inv: margin note | slanted left lower | question marks on n'a pas supposé
% margin note written sideways
% slanted
\margin{?? On n'a pas supposé pourtant que $\varphi(X)\mathrel{|\circ|}\varphi(Y)\Longrightarrow X\mathrel{|\circ|}Y$}

% inv: text | lower third | Exemple 1)
\uline{Exemple}\quad 1) Considérons l'application
% inv: display | lower third | M to Figele M
\[
M\longrightarrow\underbrace{\unsure{\mathrm{Fig\acute{e}l}|\mathrm{Fig\acute{e}le}}(M)}_{\substack{\text{figures \uline{ensemblistes}}\\\text{élémentaires}}}\qquad X\longmapsto\unsure{\mathrm{Mulomb}|\mathrm{Muloms}}(X)
\]
% inv: text | bottom third | C'est un plongement
C'est un \uline{plongement} pour $\leq$, $\ll$, $\gtrless$, qu'en est-il pour $\mathrel{|\circ|}$ ?

% inv: text | bottom | On a
On a
% inv: display | bottom | Mulomb supports
\[
\begin{aligned}
|\mathrm{Mulomb}(X)|&=\mathrm{Omb}(X)=M_{\ll X}\\
|\mathrm{Mulomb}(X)|^\circ&=\mathrm{Omb}^\circ(X)=M_{\overset{\circ}{\ll}X}
\end{aligned}
\]

\scanpage{39}
\mspage{30}
% archive sheet number: 39
% inv: text | top | donc
 donc
% inv: display | top | deleted Mulomb equivalence
\[
\deleted{\unsure{\mathrm{Mulomb}|\mathrm{Muloms}}(X)\mathrel{|\circ|}\mathrm{Mulomb}(Y)\Longleftrightarrow}
\]
% inv: display | upper third | Mulomb disjointness equivalences
\[
\begin{aligned}
\mathrm{Mulomb}(X)\mathrel{|\circ|}\mathrm{Mulomb}(Y)&\Longleftrightarrow\deleted{\unsure{M_{\overset{\circ}{\ll}X}}}\mathrm{Omb}^\circ(X)\cap\mathrm{Omb}^\circ(Y)=\varnothing\\
&\Longleftrightarrow\nexists Z\in M,\text{ avec }Z\overset{\circ}{\ll}X,\ Z\overset{\circ}{\ll}Y.
\end{aligned}
\]
% inv: deletion | upper third | i.e. ceci
\deleted{i.e. ceci :}

% inv: deletion | upper middle | Est-ce
\deleted{Est-ce}\deleted{\unsure{que}}
% inv: text | upper middle | On a donc bien
On a donc bien
% inv: display | middle | X disjoint Y implies
\[
X\mathrel{|\circ|}Y\Longrightarrow\mathrm{Mulomb}(X)\mathrel{|\circ|}\mathrm{Mulomb}(Y)
\]
% inv: insertion | below implication | i.e. X Y pas minoré
\inserted{i.e. $\{X,Y\}$ pas minoré \unsure{ds|dans} $M,\overset{\circ}{\ll}$.}

% inv: text | middle | Car si on avait
Car si on avait $Z\overset{\circ}{\ll}X$, $Z\overset{\circ}{\ll}Y$, on aurait $Z\deleted{\overset{\circ}{\ll}}\inserted{\mathrel{|\circ|}}Z$ par Mag 3 b), absurde par Mag 3 a).

% inv: text | middle lower | Mais l'inverse
Mais l'inverse est-il tjrs vrai, i.e. $X\mathrel{|\circ|}Y$ peut-il s'exprimer en termes de la relation d'ordre $\overset{\circ}{\ll}$ ?

% inv: deletion | lower middle | Ça ne marche que
\deleted{Ça \unsure{ne marche que}} \deleted{C'est OK dans}

% inv: text | lower third | Mais notons que
Mais notons que si $x,y\in L$, $x\ne y$, alors $\{x,y\}$ n'est pas minoré dans $M$ pour $\overset{\circ}{\ll}$, mais on n'aura pas pour autant, en général, $x\mathrel{|\circ|}y$\deleted{\unsure{?}} i.e. $x\parallel y$ ! (Même, dans le cas des magasins \unsure{dérivables|divisibles}.)

% inv: text | bottom | Donc le fait
Donc le fait que $\mathrel{|\circ|}$ soit induit par $\overset{\circ}{\ll}$ s'apparente à une hypothèse de nature ``quasi-ensembliste'', mais sans référence particulière aux lieux.

\scanpage{40}
\mspage{31}
% archive sheet number: 40
% inv: text | top | On peut dès lors
On peut dès lors se demander quels axiomes à mettre sur $(M,\leq,\overset{\circ}{\ll})$, pour que la relation
% inv: display | upper third | star disjointness definition
\[
X\mathrel{|\circ|}Y\Longleftrightarrow\{X,Y\}\text{ non minoré \unsure{ds|dans} }M_{\overset{\circ}{\ll}}\tag*{(*)}
\]
% inv: text | upper middle | donne avec
 donne, avec $\leq$, $\overset{\circ}{\ll}$, un magasin.

% inv: text | upper middle | Les axiomes
Les axiomes Mag 1, Mag 2 (p. 22\unsure{?}, à \unsure{un .|un \textquotesingle} \unsure{près} avec les $'$) sont clairs, \deleted{Mag $3'$ \unsure{est}} sur $\ll$, $\overset{\circ}{\ll}$, \deleted{\unsure{\&}} Mag $3'$ est automatique. Mag $4'$ \deleted{\unsure{donne une condition supplémentaire}} \inserted{résulte de Mag $2'$ a !}

% inv: deletion | middle | Autre
\deleted{\unsure{Autre}}

% inv: text | middle | Proposition
\uline{Proposition}\quad Les magasins $(M,\leq,\overset{\circ}{\ll})$ (avec $\mathrel{|\circ|}$ définissable par \unsure{(*)} ci-dessus) sont caractérisés par les conditions ;

% inv: text | middle | Mag 1 prime
\uline{Mag $1'$}\quad $\leq$, $\overset{\circ}{\ll}$ des relations d'ordre

% inv: text | middle | Mag 2 prime
\uline{Mag $2'$}\quad Comme l'homologue p. 22.

% inv: deletion | middle lower | Mag 3 prime
\deleted{\uline{Mag $3'$}}\quad\deleted{$\unsure{\forall X,Y,\ X\overset{\circ}{\ll}Y}$}

% inv: text | lower middle | C'est pourtant aussi simple
C'est \unsure{pourtant} aussi simple que le cas ensembliste\inserted{(p. 9)} ! Il n'y a pas à se mettre en frais d'un \unsure{Magquens}, comme page 24. Finalement, on a une hiérarchie de structures de plus en plus particulières, des \unsure{noti-magasins|anti-magasins} :

% inv: diagram | lower third | hierarchy of magasins
\diagram{\scriptsize\begin{tikzcd}[ampersand replacement=\&,column sep=small,row sep=small]
\& {}\arrow[dd,phantom,"{\left\{\vphantom{\begin{matrix}a\\b\\c\end{matrix}}\right.}"{anchor=center}] \& \text{Magasins généraux} \& \leq,\overset{\circ}{\ll},\mathrel{|\circ|}\quad\substack{\text{Mag 1--Mag 4}\\\text{(p. 22)}}\\
\substack{M\text{ est l'ens}\\\text{des bornes}\\\text{\unsure{principales|princip.}}} \& \& \text{Magasins }\leq,\overset{\circ}{\ll}\arrow[u,Rightarrow] \& \substack{\text{Mag }1',\text{Mag }2'=\text{Mag 2}\\\text{(cette page)}}\\
\& {} \& \text{Magasins quasi-ensemblistes}\arrow[u,Rightarrow] \& \text{Mag }1',\text{Mag }2',\text{Mag }3'\quad\text{(p. 24)}\\
\substack{L\text{ est l'ens des}\\\text{bornes princip.}} \& \{ \& \text{Magasins ensemblistes}\arrow[u,Rightarrow] \& \unsure{\mathrm{Magens}}1,\mathrm{Magens}2\quad\text{(p. 9)}
\end{tikzcd}}

% inv: text | bottom | On sera amené aussi
On sera amené aussi à regarder les magasins $\leq$, $\mathrel{|\circ|}$, qui auront d'autres avantages (qui se trouvent entre les magasins généraux, et les ensemblistes\ldots).

\scanpage{41}
\mspage{32}
% archive sheet number: 41
% inv: text | top | Exemple 2
\uline{Exemple 2}\quad Considérons l'application
% inv: display | upper third | varphi M to Figele L
\[
\varphi:M\longrightarrow\unsure{\mathrm{Fig\acute{e}l}|\mathrm{Fig\acute{e}le}}(L)\qquad X\longrightarrow\unsure{\mathrm{mulomb}|\mathrm{muloms}}(\unsure{X})
\]
% inv: deletion | upper third | On sait que c'est
\deleted{On sait que c'est}
% inv: text | upper third | Ce n'est défini
Ce n'est défini \inserted{(comme appli. à val. \unsure{ds|dans} \uuline{Figél})} que si on suppose qu'on a Mag L 1, i.e. $\forall X\in M$, $\mathrm{omb}(X)^\circ\ne\varnothing$.

% inv: text | upper middle | Dans ce cas
Dans ce cas, on a \deleted{bien}
% inv: display | upper middle | cancelled implication
\[
\deleted{X\mathrel{|\circ|}Y\Longrightarrow\mathrm{omb}(X)}
\]
% inv: display | middle | mulomb supports
\[
\begin{aligned}
|\mathrm{mulomb}(X)|&=\mathrm{omb}(X)\\
|\mathrm{mulomb}(X)|^\circ&=\mathrm{omb}(X)^\circ
\end{aligned}
\]
% inv: text | middle | et
et
% inv: display | middle | disjoint mulomb equivalence
\[
\mathrm{mulomb}(X)\mathrel{|\circ|}\mathrm{mulomb}(Y)\Longleftrightarrow\mathrm{omb}(X)^\circ\cap\mathrm{omb}(Y)^\circ=\unsure{\varnothing}
\]
% inv: text | middle lower | Donc on a bien
Donc on a bien
% inv: display | middle lower | implication to mulomb Y
\[
X\mathrel{|\circ|}Y\Longrightarrow\mathrm{mulomb}(X)\mathrel{|\circ|}\mathrm{mulomb}\,Y.
\]

% inv: deletion | lower middle | Cela nous montre
\deleted{Cela nous \unsure{montre}}
% inv: text | lower third | Ces deux exemples
Ces deux exemples nous montrent en tout cas que la notion de comorphisme proposée à la page 27, n'est pas trop idiote. (même s'il convient sans doute de la généraliser). Troisième exemple :

% inv: text | bottom third | Exemple 3
\uline{Exemple 3}\quad $f:\unsure{E|\mathbb{E}}'\longrightarrow E$ application ensembliste surjective. Elle induit un comorphisme injectif des magasins ensemblistes associés
% inv: display | bottom | Figel E to Figel E prime
\[
\mathrm{Figel}(E)\longrightarrow\mathrm{Figel}(E')
\]
